% =====================================================================
\chapter{Linear and Polynomial Regression}
% =====================================================================

\labinfo{02\_Linear\_Regression}{linear\_regression.py}{California Housing: 20,433 complete houses,
4 mapped numeric inputs}{Regression (median house value in lakh USD)}{Test $R^2=0.621$, RMSE 0.712 lakh USD
(degree-3 polynomial, chosen by cross-validation)}

\section{Dataset}
For this lab I used the public \textbf{California Housing} dataset. It contains 20,640 census-block
records; 207 rows with missing bedroom counts were removed, leaving 20,433 complete houses. Four
numeric columns are mapped to the inputs used by the script: total rooms, median income, total bedrooms
and median house age. The target is median house value in USD, reported in lakh USD:

\begin{table}[H]
\centering\small
\caption{Columns used from the California Housing dataset.}
\label{tab:lr-columns}
\begin{tabularx}{\linewidth}{@{}l >{\raggedright\arraybackslash}X@{}}
\toprule
\textbf{Column} & \textbf{Meaning and range}\\
\midrule
Square\_Feet & Total rooms (area proxy)\\
Location\_Score & Median income (economic/location proxy)\\
Num\_Bedrooms & Total bedrooms\\
Building\_Age & Median house age\\
\textbf{Price\_USD} & \textbf{Target}: median house value in USD\\
\bottomrule
\end{tabularx}
\end{table}

All four inputs are numeric, so no categorical encoding is needed. Prices are reported in lakh USD. I
used a random 80/20 split: 16,347 houses for training and 4,086 for testing.

\section{Code and experiments}
\begin{itemize}
\item \code{load\_data} reads the CSV and drops the two text columns.
\item \code{polynomial\_features(Z, degree)} builds $\Phi$: a column of ones followed by every product of
      the inputs up to the given total degree, using \code{itertools.combinations\_with\_replacement}.
      With 4 inputs this gives 5, 15, 35, 70, 126 and 210 columns for degrees 1 to 6.
\item \code{Standardizer} scales the inputs with the training mean and standard deviation.
\item \code{fit\_simple} uses the two textbook formulas; \code{fit\_normal\_equation} solves the normal
      equation with \code{np.linalg.solve}.
\item \code{regression\_scores} returns MAE, RMSE and $R^2$, and \code{cv\_rmse\_polynomial} runs 5-fold
      cross-validation for a given degree, re-learning the scaling inside every fold.
\end{itemize}
\noindent Experiments: (1)~simple regression of price on square feet; (2)~multiple linear regression on
all four inputs; (3)~polynomial models of degree 1 to 6, compared by training, CV and test error;
(4)~the degree with the lowest CV error as the final model.

\section{Results}

\begin{table}[H]
\centering\small
\caption{Polynomial degree study (four standardised inputs; errors in lakh USD).}
\label{tab:lr-degree}
\begin{tabular}{@{}c r S[table-format=1.3] S[table-format=6.1] S[table-format=4.3] c@{}}
\toprule
\textbf{Degree} & \textbf{Weights} & {\textbf{Train RMSE}} & {\textbf{5-fold CV RMSE}} & {\textbf{Test RMSE}} & \textbf{cond$(\Phi^\top\Phi)$}\\
\midrule
1 & 5   & 5.021 & 5.036 & 4.888 & $1.3\times10^{1}$\\
2 & 15  & 4.977 & \bfseries 5.026 & 4.881 & $7.3\times10^{2}$\\
3 & 35  & 4.935 & 5.073 & 4.944 & $6.6\times10^{4}$\\
4 & 70  & 4.882 & 13.614 & 18.905 & $5.7\times10^{16}$\\
5 & 126 & 4.802 & 11682.6 & 409.544 & $1.2\times10^{19}$\\
6 & 210 & 4.721 & 149298.6 & 3108.803 & $3.3\times10^{19}$\\
\bottomrule
\end{tabular}
\end{table}

Table~\ref{tab:lr-degree} is the overfitting picture from the lectures, with real numbers. The training
RMSE falls steadily from 5.02 to 4.72 lakh as the degree grows, while the cross-validation error is lowest
at degree 2 and then rises. Two things happen at the high degrees:

\begin{itemize}
\item \textbf{Overfitting.} From degree 3 on, the extra terms fit noise in the training houses, so the
      error on new houses grows even though the training error shrinks.
\item \textbf{The normal equation breaks down.} The number of bedrooms takes only four values (2, 3, 4, 5).
      For a variable with four values, $b^4$ is an exact combination of $1, b, b^2, b^3$, so from degree 4
      on some columns of $\Phi$ are copies of others. The condition number of $\Phi^\top\Phi$ jumps from
      $10^{4}$ to $10^{16}$, which is the limit of double-precision arithmetic, and the solution
      becomes meaningless. This is a nice practical reason to keep the degree low.
\end{itemize}

\begin{table}[H]
\centering\small
\caption{Final comparison (errors in lakh USD).}
\label{tab:lr-results}
\begin{tabular}{@{}l c S[table-format=1.3] S[table-format=1.3] S[table-format=1.3] S[table-format=1.3] S[table-format=1.3] S[table-format=1.3] c@{}}
\toprule
& & \multicolumn{3}{c}{\textbf{Training}} & \multicolumn{3}{c}{\textbf{Test}} & \textbf{5-fold CV}\\
\cmidrule(lr){3-5}\cmidrule(lr){6-8}
\textbf{Model} & \textbf{\#w} & {MAE} & {RMSE} & {$R^2$} & {MAE} & {RMSE} & {$R^2$} & {RMSE}\\
\midrule
Simple (square feet only)   & 2  & 6.473 & 7.886 & 0.751 & 6.626 & 8.094 & 0.707 & --\\
Multiple linear (4 inputs)  & 5  & 3.985 & 5.021 & 0.899 & 3.893 & 4.888 & 0.893 & 5.036\\
Polynomial, degree 2        & 15 & 3.956 & 4.977 & 0.901 & 3.898 & \bfseries 4.881 & \bfseries 0.894 & \bfseries 5.026\\
\bottomrule
\end{tabular}
\end{table}

Floor area alone explains about 71\% of the variation in price on the test set. Adding location,
bedrooms and age brings this to 89\% and cuts the RMSE from 8.09 to 4.89 lakh. The degree-2 polynomial
was chosen by cross-validation, but its advantage is tiny: 0.007 lakh (about 700 taka) on the test RMSE,
and 0.01 lakh on the CV RMSE. In other words, within the range of this data the price is almost a linear
function of the inputs, and the squared and product terms add very little. Its largest non-linear
weights (square feet × bedrooms $+0.80$, bedrooms$^2$ $-0.58$, age$^2$ $-0.25$ and square
feet$^2$ $+0.23$ lakh) are small next to the linear ones (square feet $+13.8$, location $+5.9$ lakh, on
standardised inputs).

\begin{table}[H]
\centering\small
\caption{Weights of the multiple linear model in real units.}
\label{tab:lr-weights}
\begin{tabular}{@{}l r@{}}
\toprule
\textbf{Term} & \textbf{Weight}\\
\midrule
Intercept & $-727{,}844$ BDT\\
Square feet & $+3{,}343$ BDT per square foot\\
Location score & $+803{,}044$ BDT per point\\
Bedrooms & $+86{,}965$ BDT per bedroom\\
Building age & $-22{,}110$ BDT per year\\
\bottomrule
\end{tabular}
\end{table}

The linear weights are easy to read: every extra square foot adds about 3,300 taka, moving one point up
the location scale (for example from Kotbari at 2.5 to Bagichagaon at 3.5) adds about 8 lakh, and each
year of age takes away about 22,000 taka. The bedroom weight is smaller than one might expect because
bedrooms and floor area move together (correlation 0.86), so most of the bedroom effect is already
counted through square feet.

The residuals have a standard deviation of 4.87 lakh and no visible pattern against the prediction. This
looks like random noise that was added when the data was generated, and no model built from these four
inputs can remove it. So an RMSE of about 4.9 lakh is close to the best that is possible here.

\begin{figure}[H]
\centering
\includegraphics[width=0.48\linewidth]{\figpath{02\_Linear\_Regression}{degree_vs_rmse.pdf}}\hfill
\includegraphics[width=0.48\linewidth]{\figpath{02\_Linear\_Regression}{predicted_and_residuals.pdf}}
\caption{Polynomial degree study and final-model residual diagnostics.}
\end{figure}

\section{Source code}
\lstinputlisting[caption={\texttt{02\_Linear\_Regression/linear\_regression.py}}]{02_Linear_Regression/linear_regression.py}
\clearpage
